\documentclass[]{book}
\usepackage{lmodern}
\usepackage{amssymb,amsmath}
\usepackage{ifxetex,ifluatex}
\usepackage{fixltx2e} % provides \textsubscript
\ifnum 0\ifxetex 1\fi\ifluatex 1\fi=0 % if pdftex
  \usepackage[T1]{fontenc}
  \usepackage[utf8]{inputenc}
\else % if luatex or xelatex
  \ifxetex
    \usepackage{mathspec}
  \else
    \usepackage{fontspec}
  \fi
  \defaultfontfeatures{Ligatures=TeX,Scale=MatchLowercase}
\fi
% use upquote if available, for straight quotes in verbatim environments
\IfFileExists{upquote.sty}{\usepackage{upquote}}{}
% use microtype if available
\IfFileExists{microtype.sty}{%
\usepackage{microtype}
\UseMicrotypeSet[protrusion]{basicmath} % disable protrusion for tt fonts
}{}
\usepackage[margin=1.5in]{geometry}
\usepackage{hyperref}
\hypersetup{unicode=true,
            pdftitle={Template},
            pdfborder={0 0 0},
            breaklinks=true}
\urlstyle{same}  % don't use monospace font for urls
\usepackage{natbib}
\bibliographystyle{apalike}
\usepackage{color}
\usepackage{fancyvrb}
\newcommand{\VerbBar}{|}
\newcommand{\VERB}{\Verb[commandchars=\\\{\}]}
\DefineVerbatimEnvironment{Highlighting}{Verbatim}{commandchars=\\\{\}}
% Add ',fontsize=\small' for more characters per line
\usepackage{framed}
\definecolor{shadecolor}{RGB}{248,248,248}
\newenvironment{Shaded}{\begin{snugshade}}{\end{snugshade}}
\newcommand{\KeywordTok}[1]{\textcolor[rgb]{0.13,0.29,0.53}{\textbf{#1}}}
\newcommand{\DataTypeTok}[1]{\textcolor[rgb]{0.13,0.29,0.53}{#1}}
\newcommand{\DecValTok}[1]{\textcolor[rgb]{0.00,0.00,0.81}{#1}}
\newcommand{\BaseNTok}[1]{\textcolor[rgb]{0.00,0.00,0.81}{#1}}
\newcommand{\FloatTok}[1]{\textcolor[rgb]{0.00,0.00,0.81}{#1}}
\newcommand{\ConstantTok}[1]{\textcolor[rgb]{0.00,0.00,0.00}{#1}}
\newcommand{\CharTok}[1]{\textcolor[rgb]{0.31,0.60,0.02}{#1}}
\newcommand{\SpecialCharTok}[1]{\textcolor[rgb]{0.00,0.00,0.00}{#1}}
\newcommand{\StringTok}[1]{\textcolor[rgb]{0.31,0.60,0.02}{#1}}
\newcommand{\VerbatimStringTok}[1]{\textcolor[rgb]{0.31,0.60,0.02}{#1}}
\newcommand{\SpecialStringTok}[1]{\textcolor[rgb]{0.31,0.60,0.02}{#1}}
\newcommand{\ImportTok}[1]{#1}
\newcommand{\CommentTok}[1]{\textcolor[rgb]{0.56,0.35,0.01}{\textit{#1}}}
\newcommand{\DocumentationTok}[1]{\textcolor[rgb]{0.56,0.35,0.01}{\textbf{\textit{#1}}}}
\newcommand{\AnnotationTok}[1]{\textcolor[rgb]{0.56,0.35,0.01}{\textbf{\textit{#1}}}}
\newcommand{\CommentVarTok}[1]{\textcolor[rgb]{0.56,0.35,0.01}{\textbf{\textit{#1}}}}
\newcommand{\OtherTok}[1]{\textcolor[rgb]{0.56,0.35,0.01}{#1}}
\newcommand{\FunctionTok}[1]{\textcolor[rgb]{0.00,0.00,0.00}{#1}}
\newcommand{\VariableTok}[1]{\textcolor[rgb]{0.00,0.00,0.00}{#1}}
\newcommand{\ControlFlowTok}[1]{\textcolor[rgb]{0.13,0.29,0.53}{\textbf{#1}}}
\newcommand{\OperatorTok}[1]{\textcolor[rgb]{0.81,0.36,0.00}{\textbf{#1}}}
\newcommand{\BuiltInTok}[1]{#1}
\newcommand{\ExtensionTok}[1]{#1}
\newcommand{\PreprocessorTok}[1]{\textcolor[rgb]{0.56,0.35,0.01}{\textit{#1}}}
\newcommand{\AttributeTok}[1]{\textcolor[rgb]{0.77,0.63,0.00}{#1}}
\newcommand{\RegionMarkerTok}[1]{#1}
\newcommand{\InformationTok}[1]{\textcolor[rgb]{0.56,0.35,0.01}{\textbf{\textit{#1}}}}
\newcommand{\WarningTok}[1]{\textcolor[rgb]{0.56,0.35,0.01}{\textbf{\textit{#1}}}}
\newcommand{\AlertTok}[1]{\textcolor[rgb]{0.94,0.16,0.16}{#1}}
\newcommand{\ErrorTok}[1]{\textcolor[rgb]{0.64,0.00,0.00}{\textbf{#1}}}
\newcommand{\NormalTok}[1]{#1}
\usepackage{longtable,booktabs}
\usepackage{graphicx,grffile}
\makeatletter
\def\maxwidth{\ifdim\Gin@nat@width>\linewidth\linewidth\else\Gin@nat@width\fi}
\def\maxheight{\ifdim\Gin@nat@height>\textheight\textheight\else\Gin@nat@height\fi}
\makeatother
% Scale images if necessary, so that they will not overflow the page
% margins by default, and it is still possible to overwrite the defaults
% using explicit options in \includegraphics[width, height, ...]{}
\setkeys{Gin}{width=\maxwidth,height=\maxheight,keepaspectratio}
\IfFileExists{parskip.sty}{%
\usepackage{parskip}
}{% else
\setlength{\parindent}{0pt}
\setlength{\parskip}{6pt plus 2pt minus 1pt}
}
\setlength{\emergencystretch}{3em}  % prevent overfull lines
\providecommand{\tightlist}{%
  \setlength{\itemsep}{0pt}\setlength{\parskip}{0pt}}
\setcounter{secnumdepth}{5}
% Redefines (sub)paragraphs to behave more like sections
\ifx\paragraph\undefined\else
\let\oldparagraph\paragraph
\renewcommand{\paragraph}[1]{\oldparagraph{#1}\mbox{}}
\fi
\ifx\subparagraph\undefined\else
\let\oldsubparagraph\subparagraph
\renewcommand{\subparagraph}[1]{\oldsubparagraph{#1}\mbox{}}
\fi

%%% Use protect on footnotes to avoid problems with footnotes in titles
\let\rmarkdownfootnote\footnote%
\def\footnote{\protect\rmarkdownfootnote}

%%% Change title format to be more compact
\usepackage{titling}

% Create subtitle command for use in maketitle
\providecommand{\subtitle}[1]{
  \posttitle{
    \begin{center}\large#1\end{center}
    }
}

\setlength{\droptitle}{-2em}

  \title{Template}
    \pretitle{\vspace{\droptitle}\centering\huge}
  \posttitle{\par}
    \author{}
    \preauthor{}\postauthor{}
      \predate{\centering\large\emph}
  \postdate{\par}
    \date{August 2019}

\usepackage{booktabs}

\usepackage{epsfig}
\usepackage{epstopdf}
\usepackage{rotate}
\usepackage{graphicx}
\usepackage{hyperref}
\usepackage{alphalph}
\usepackage{caption}
\usepackage[hang,flushmargin]{footmisc}
\usepackage{framed}
\usepackage{xcolor}
\usepackage{verbatim} 

\usepackage{bm}
\setcounter{MaxMatrixCols}{20}
\newcommand{\Var}{\mathrm{Var}}
\newcommand{\SD}{\mathrm{SD}}
\newcommand{\Cov}{\mathrm{Cov}}
\newcommand{\fx}{f({\bf x})}
\newcommand\R{{\textsf R~}}
\newcommand\Rst{\textsf{RStudio}}

% spacing between environments
\usepackage{amsthm}
\makeatletter
\def\thm@space@setup{%
  \thm@preskip=15pt plus 2pt minus 4pt
  \thm@postskip=\thm@preskip
}
\makeatother


% Title format
\usepackage{titling}
\pretitle{\Huge\sffamily}
\posttitle{\par\vskip 0.5em}
\predate{\LARGE\sffamily}
\postdate{\par}

\urlstyle{tt}

\begin{document}
\maketitle

{
\setcounter{tocdepth}{1}
\tableofcontents
}
\hypertarget{introduction}{%
\chapter*{Introduction}\label{introduction}}
\addcontentsline{toc}{chapter}{Introduction}

\hypertarget{table-of-contents}{%
\section*{Table of Contents}\label{table-of-contents}}
\addcontentsline{toc}{section}{Table of Contents}

Here, we outline how the guide is organized into parts.
\textbf{TEST}

****TESTS TESTS TESTS*******
****MYTESTS MYTESTS MYTESTS*******
1. First, we\ldots{}
2. Second, we\ldots{}
3. Lastly, we\ldots{}

\href{https://www.rstudio.com/wp-content/uploads/2015/03/devtools-cheatsheet.pdf}{Here we provide an outside link to important content} which puts some useful information for this tutorial/workshop at our fingertips.

Here we specify where people can provide feedback! Please email \href{mailto:help@iq.harvard.edu}{\nolinkurl{help@iq.harvard.edu}}

\hypertarget{authors-and-sources}{%
\section*{Authors and Sources}\label{authors-and-sources}}
\addcontentsline{toc}{section}{Authors and Sources}

Here we acknowledge a few people who helped make this tutorial/workshop possible. We also reference any sources that material was taken from.

\hypertarget{this-is-a-section}{%
\chapter{This is a section}\label{this-is-a-section}}

******THIS IS A TEST*********

!!!!!!!THIS IS ANOTHER TEST!!!!!!!!!!!

Here's how we embed an image:

\includegraphics{images/readme-license.png}

And here's how we embed a URL \href{https://dss.iq.harvard.edu/}{DSS}.

\hypertarget{this-is-a-subsection}{%
\section{This is a subsection}\label{this-is-a-subsection}}

Text for my subsection.

\hypertarget{this-is-a-subsubsection}{%
\subsection{This is a subsubsection}\label{this-is-a-subsubsection}}

Text for my subsubsection.

\hypertarget{lists}{%
\section{Lists}\label{lists}}

This is a numbered list:

\begin{enumerate}
\def\labelenumi{\arabic{enumi}.}
\tightlist
\item
  Item 1
\end{enumerate}

\begin{itemize}
\tightlist
\item
  sub-item1
\item
  sub-item2
\end{itemize}

\begin{enumerate}
\def\labelenumi{\arabic{enumi}.}
\setcounter{enumi}{1}
\tightlist
\item
  Item 2
\end{enumerate}

\begin{itemize}
\tightlist
\item
  sub-item1
\item
  sub-item2
\end{itemize}

This is an unnumbered list:

\begin{itemize}
\tightlist
\item
  Item 1

  \begin{itemize}
  \tightlist
  \item
    sub-item1
  \item
    sub-item2
  \end{itemize}
\item
  Item 2

  \begin{itemize}
  \tightlist
  \item
    sub-item1
  \item
    sub-item2
  \end{itemize}
\end{itemize}

\hypertarget{text-emphasis}{%
\section{Text emphasis}\label{text-emphasis}}

We can emphasize text by using \textbf{bold} or \emph{italics}

\hypertarget{embedding-code}{%
\section{Embedding code}\label{embedding-code}}

We can embed R code in blocks. The blocks can either be run (evaluated) or just printed:

\begin{Shaded}
\begin{Highlighting}[]
\CommentTok{# install.packages('reticulate')}
\KeywordTok{library}\NormalTok{(reticulate)}
\KeywordTok{use_python}\NormalTok{(}\StringTok{"~/anaconda3/bin/python"}\NormalTok{) }\CommentTok{# use anaconda python on Mac}
\KeywordTok{use_python}\NormalTok{(}\StringTok{"C:/Users/<username>/Anaconda3/python.exe"}\NormalTok{) }\CommentTok{# use anaconda python on Windows}
\end{Highlighting}
\end{Shaded}

We can also \href{https://rstudio.github.io/reticulate/articles/r_markdown.html}{embed Python code blocks within Rmarkdown} using the \texttt{reticulate} package, once it is \href{https://rstudio.github.io/reticulate/index.html}{installed}.

\begin{Shaded}
\begin{Highlighting}[]
\NormalTok{flights }\OperatorTok{=}\NormalTok{ pandas.read.csv(}\StringTok{'flights.csv'}\NormalTok{)}
\NormalTok{flights }\OperatorTok{=}\NormalTok{ flights[flights[}\StringTok{'dest'}\NormalTok{] }\OperatorTok{==} \StringTok{'ORD'}\NormalTok{]}
\NormalTok{flights }\OperatorTok{=}\NormalTok{ flights.dropna()}
\end{Highlighting}
\end{Shaded}

Or we can format code inline (not evaluated): \texttt{print()}


\end{document}
